% !TITLE 锦瑟
% !AUTHOR 李商隐
% !DYNASTY 唐
% !GENRE 七言律诗
% !ORDER 54

\begin{poem}
锦瑟无端五十弦，一弦一柱思华年。
庄生晓梦迷蝴蝶，望帝春心托杜鹃。
沧海月明珠有泪，蓝田日暖玉生烟。
此情可待成追忆？只是当时已惘然。
\end{poem}

\begin{appreciation}
《锦瑟》是李商隐最难解的一首诗，也是他最动人的一首诗。千百年来，无数人试图给它一个确定的解释——是悼亡？是自伤？是政治寄托？但也许这首诗的魅力，恰恰在于它拒绝被解释。

锦瑟无端五十弦——瑟是一种弦乐器，通常二十五弦，这里说五十弦，是用了典故：传说中泰帝使素女鼓五十弦瑟，音悲，帝禁不止，故破其瑟为二十五弦。无端是无缘无故的意思。诗人看着锦瑟，忽然觉得那五十根弦毫无来由地摆在那里，每一根都牵出一缕思绪。一弦一柱思华年——每一弦、每一柱，都让他想起逝去的年华。华年就是美好的年华，这里指已经消逝的青春岁月。

庄生晓梦迷蝴蝶——庄子梦见自己变成了蝴蝶，醒来后不知是自己梦见了蝴蝶，还是蝴蝶梦见了自己。这个典故说的是物我两忘、真实与虚幻的界限。望帝春心托杜鹃——望帝是古蜀国的君主，死后化为杜鹃鸟，每到春天便啼血哀鸣。这个典故说的是无尽的哀怨与执着。李商隐把这两个典故并置，不做任何解释，像两面镜子互相映照，让人在其中看见自己的迷茫与哀伤。

沧海月明珠有泪——月光下，沧海中的珍珠像眼泪一样闪烁。蓝田日暖玉生烟——阳光照耀下，蓝田的玉石升起朦胧的烟霭。这两句对仗极为工整，意象极为优美，却又极为虚幻。珠泪是冷的，玉烟是暖的；沧海是远的，蓝田是近的。诗人似乎在写某种可望而不可即的东西——也许是理想，也许是爱情，也许是人生本身。

此情可待成追忆？只是当时已惘然——这份情感何必等到将来再去追忆呢？其实在当时就已经迷惘了。可待是岂待的意思，惘然是迷惘的样子。这不是一个问句，而是一个陈述：人生本来就是惘然的，即使身处其中，也已知道终将失去。

《锦瑟》的妙处，在于它像一首无题的音乐，你听不懂每一个音符的含义，却能感受到整体的悲伤。李商隐用四个典故、两幅画面，搭建了一座迷宫，而他自己，也迷失在其中。
\end{appreciation}
